\subsection{Sprint 1}

Para esta sprint estavam previstas atividades como modelagem do banco de dados,
criação dos scripts de tratamento de dados, mapeamento do SQLAlchemy,
correções do Figma conforme demanda do cliente e uma spike para analisar como
coletar e quais dados seriam possíveis de serem coletados. Da minha parte, como
AGES II, era esperada uma grande participação nas questões relacionadas ao banco
de dados, bem como colaboração com colegas, resultando em uma participação
indireta nas outras tarefas.

Quanto à parte do banco de dados, a equipe teve uma boa quantidade de conversas
que ajudaram a avançar a modelagem, permitindo a criação de uma estrutura
definitiva relativa aos dados que tínhamos até o momento, sempre sujeita a
alterações dependendo do resultado das pesquisas.

Os scripts de tratamento já haviam sido discutidos e estavam em processo de
implementação. A implementação do SQLAlchemy, por sua vez, caminhava lado a
lado com a modelagem. As mudanças do Figma foram realizadas de maneira ágil e
estavam prontas para entrega desde o começo da sprint. Já a spike, apesar de ter
resultado em boas descobertas, mostrou que o processo de coleta não era tão fácil
quanto parecia e demandaria mais tempo de pesquisa. Ainda assim, foi um
progresso interessante ter uma boa ideia de quais ferramentas seriam utilizadas.

Da minha parte, participei ativamente de discussões relativas à modelagem do
banco de dados e ao tratamento dos dados, tanto em sala de aula quanto fora dela.
Além disso, colaborei com meus colegas na implementação das entidades e dos
relacionamentos no SQLAlchemy. Quanto à spike, participei com menor intensidade
do que no banco de dados, mas cheguei a pesquisar a Graph \ac{api} da Meta e a
viabilidade de uso para Instagram e Facebook. No entanto, tive dificuldade ao ler a
documentação, devido à terminologia mais complexa utilizada para definir a
frequência máxima de requisições possíveis para a \ac{api}.

Em um projeto diferente dos normais da AGES, onde pesquisa e banco de dados
são dois dos principais fatores, a definição padrão de uma sprint acabou deixando
a desejar. Ainda assim, acreditávamos que não seria tão problemático estarmos
"atrasados" em comparação com um fluxo normal da AGES, visto que, após
encaminhar banco de dados e métodos de coleta, o desenvolvimento do aplicativo
se tornaria mais direto. Além disso, o fato de o front-end não demandar tantos
recursos humanos possibilitaria agilizar as entregas nas próximas sprints e
terminar no tempo previsto.

Um projeto como este ensina a pensar fora da caixa e a estar aberto para novas
funcionalidades, métodos e agendas. A peculiaridade do projeto trouxe certo receio,
principalmente devido ao peso desta sprint, tendo em mente que o banco de dados,
que inicialmente parecia algo trivial, provou ser uma grande barreira por causa da
finalidade do projeto: o processamento de uma base massiva de dados. Também
foi possível aprender com os acertos. Um ponto positivo observado no time foi a
união, com diálogo regular entre a maioria dos colegas, criando um ambiente de
trabalho saudável e composto por pessoas que mantinham comunicação inclusive
fora da AGES. Esse ponto foi positivo para o desenvolvimento do projeto, para a
sincronia da equipe e para o desenvolvimento pessoal, pois o compartilhamento de
experiências acelerou a absorção de conhecimento.

Com o esqueleto do banco de dados montado e as relações já implementadas no
SQLAlchemy, acredito que havia três caminhos futuros para o time. Um deles era a
continuação da spike de pesquisa, com o intuito de chegar a resultados mais
concretos sobre a coleta de dados possivelmente relacionáveis a um score de
presença digital. Os outros dois caminhos eram relativos à programação, um no
front-end e outro no back-end. Da minha parte, estava mais inclinado a trabalhar
com o back-end, já que fiquei no front-end durante a AGES passada e estudei mais
as tecnologias de back-end do que de front-end. Ainda assim, me mantive
disponível para ajudar o time onde fosse necessário.
